%!TEX root = ../memoria.tex
% =============================================================================
%  APÉNDICE B — Entorno y dependencias del proyecto
% =============================================================================

El código fuente completo de la aplicación se encuentra disponible en un
repositorio público, accesible en la siguiente dirección:

\begin{center}
\url{https://github.com/yagogl/conectamayor-TFG}
\end{center}

El repositorio incluye, además del código, un fichero \texttt{README} con
las instrucciones detalladas de instalación y puesta en marcha, así como
los datos de acceso de un escenario de demostración que permite explorar
la aplicación con contenido de ejemplo.

\section{Entorno de ejecución\label{SEC:APP_ENTORNO}}

La aplicación se ha desarrollado y probado sobre el siguiente entorno:

\begin{itemize}
  \item Sistema operativo: Ubuntu 24.04 LTS.
  \item Lenguaje: Python 3.12.3.
  \item Framework web: Django 6.0.4.
  \item Base de datos: SQLite, motor incluido por defecto en Django.
  \item Capa de presentación: Bootstrap 5 y JavaScript estándar.
\end{itemize}

\section{Dependencias\label{SEC:APP_DEPENDENCIAS}}

Las dependencias del proyecto se gestionan mediante el fichero
\texttt{requirements.txt}, que permite reconstruir el entorno en cualquier
sistema compatible con una única orden:

\begin{center}
\texttt{pip install -r requirements.txt}
\end{center}

El contenido de dicho fichero es el siguiente:

\begin{lstlisting}[language={},basicstyle=\ttfamily\small,frame=single,numbers=none]
asgiref==3.11.1
Django==6.0.4
pillow==12.2.0
sqlparse==0.5.5
\end{lstlisting}

El reducido número de dependencias responde a una decisión de diseño del proyecto, que busca minimizar las librerías externas y apoyarse en lo que ya ofrece Django. Las instrucciones completas para clonar el repositorio, instalar las dependencias y ejecutar la aplicación, incluida la carga opcional de los datos de demostración, están en el fichero \texttt{README} del repositorio.
